\BourbakiExerciseGroup

\begin{enumerate}
\def\labelenumi{\arabic{enumi})}
\item
  求出两个元素的集合或三个元素的集合上的所有可能的拓扑。
\item
  \begin{enumerate}
  \def\labelenumii{\alph{enumii})}
  \tightlist
  \item
    设 X 是一个有序集。证明区间所成的集
  \end{enumerate}
\end{enumerate}

\[
[ x, \rightarrow [ \quad (\text { resp. } ] \leftarrow , x ])
\]

是 X 上一个拓扑的基；这一拓扑称为 X 的右拓扑（相应地，左拓扑）。在右拓扑中，开集的任意交是开集，且 \(\{x\}\) 的闭包是区间 {]}\(\leftarrow\), x{]}。

\begin{enumerate}
\def\labelenumi{\alph{enumi})}
\setcounter{enumi}{1}
\item
  若拓扑空间满足以下条件，则称它为 Kolmogoroff 空间：给定任意两个不同的点 \(x\), \(x'\)（属于 \(X\)），存在这两点之一的一个邻域，它不含另一点。证明带有右拓扑的有序集是 Kolmogoroff 空间。
\item
  设 X 是一个 Kolmogoroff 空间，其中开集的每个交都是开集。证明 \(x \in \overline{\{x'\}}\) 是 X 中 x 与 \(x'\) 之间的一个序关系，且若把这个关系记作 \(x \leqslant x'\)，则 X 上给定的拓扑与由这个排序所确定的右拓扑相同。
\item
  由 c) 推出：若 X 是 Kolmogoroff 空间，则 X 的每个有限非空子集至少有一个孤立点。因此，若 X 没有孤立点，则 X 的每个非空开子集都是无限的。
\end{enumerate}

\begin{enumerate}
\def\labelenumi{\arabic{enumi})}
\setcounter{enumi}{2}
\tightlist
\item
  若 \(\mathbf{A}\) 是拓扑空间 \(\mathbf{X}\) 的任一子集，用 \(\alpha(\mathbf{A})\) 表示 \(\overline{\mathbf{A}}\)，用 \(\beta(\mathbf{A})\) 表示 \(\overline{\mathbf{A}}\)。则 \(\mathbf{A} \subset \mathbf{B}\) 蕴含 \(\alpha(\mathbf{A}) \subset \alpha(\mathbf{B})\) 与 \(\beta(\mathbf{A}) \subset \beta(\mathbf{B})\)。\(a)\) 证明：若 \(\mathbf{A}\) 是开的，则 \(\mathbf{A} \subset \alpha(\mathbf{A})\)；若 \(\mathbf{A}\) 是闭的，则 \(\beta(\mathbf{A}) \subset \mathbf{A}\)。
\end{enumerate}

\begin{enumerate}
\def\labelenumi{\alph{enumi})}
\setcounter{enumi}{1}
\item
  由 a) 推出：若 A 是 X 的任一子集，则 \(\alpha(\alpha(A)) = \alpha(A)\) 且 \(\beta(\beta(A)) = \beta(A)\)。
\item
  给出集合 A 的一个例子 \(*\)（在实直线上）\(_{*}\)，使得七个集合 A、\(\mathring{A}\)、\(\overline{A}\)、\(\alpha(A)\)、\(\beta(A)\)、\(\beta(\overline{A})\)、\(\alpha(\mathring{A})\) 两两不同，且除以下关系外不满足任何包含关系：\(\mathring{A} \subset A \subset \overline{A}\)，\(\mathring{A} \subset \alpha(\mathring{A}) \subset \beta(A) \subset \overline{A}\)，\(\mathring{A} \subset \alpha(A) \subset \beta(\overline{A}) \subset \overline{A}\)，\(\alpha(\mathring{A}) \subset \alpha(A)\)，\(\beta(A) \subset \beta(\overline{A})\)。
\item
  证明：若 U, V 是两个开集，使得 \(U \cap V = \varnothing\)，则也有 \(\alpha(U) \cap \alpha(V) = \varnothing\) {[}用 b){]}。
\end{enumerate}

\begin{enumerate}
\def\labelenumi{\arabic{enumi})}
\setcounter{enumi}{3}
\tightlist
\item
  \begin{enumerate}
  \def\labelenumii{\alph{enumii})}
  \tightlist
  \item
    在实直线上给出两个开集 A, B 的一个例子 *（在实直线上）*，使得四个集合 A ∩ B、B ∩ A̅、A ∩ B 与 A̅ ∩ B 两两不同。
  \end{enumerate}
\end{enumerate}

\begin{enumerate}
\def\labelenumi{\alph{enumi})}
\setcounter{enumi}{1}
\tightlist
\item
  * 在实直线上给出两个区间 A, B 的一个例子，使得 A ∩ B 不含于 A ∩ B。*
\end{enumerate}

\begin{enumerate}
\def\labelenumi{\arabic{enumi})}
\setcounter{enumi}{4}
\tightlist
\item
  若 A 是拓扑空间 X 的任一子集，A 的边界记作 \(\operatorname{Fr}(\mathbf{A})\)。
\end{enumerate}

\begin{enumerate}
\def\labelenumi{\alph{enumi})}
\item
  证明 \(\operatorname{Fr}(\overline{\mathbf{A}}) \subset \operatorname{Fr}(\mathbf{A})\)，\(\operatorname{Fr}(\mathring{\mathbf{A}}) \subset \operatorname{Fr}(\mathbf{A})\)，并给出一个例子 *（在实直线上）*，其中这三个集合两两不同。
\item
  设 A, B 是 X 的两个子集；证明
\end{enumerate}

\[
\operatorname{Fr} (\mathbf {A} \cup \mathbf {B}) \subset \operatorname{Fr} (\mathbf {A}) \cup \operatorname{Fr} (\mathbf {B})
\]

并给出一个例子 *（在实直线上）*，其中这些集合两两不同。若 \(\overline{\mathbf{A}} \cap \overline{\mathbf{B}} = \emptyset\)，证明 \(\operatorname{Fr}(\mathbf{A} \cup \mathbf{B}) = \operatorname{Fr}(\mathbf{A}) \cup \operatorname{Fr}(\mathbf{B})\)。

\begin{enumerate}
\def\labelenumi{\alph{enumi})}
\setcounter{enumi}{2}
\tightlist
\item
  若 \(\mathbf{A}\) 与 \(\mathbf{B}\) 在 \(\mathbf{X}\) 中是开的，证明
\end{enumerate}

\[
\begin{array}{r l} & (\mathbf {A} \cap \operatorname{Fr} (\mathbf {B})) \cup (\mathbf {B} \cap \operatorname{Fr} (\mathbf {A})) \subset \operatorname{Fr} (\mathbf {A} \cap \mathbf {B}) \\ & \qquad \subset (\mathbf {A} \cap \operatorname{Fr} (\mathbf {B})) \cup (\mathbf {B} \cap \operatorname{Fr} (\mathbf {A})) \cup (\operatorname{Fr} (\mathbf {A}) \cap \operatorname{Fr} (\mathbf {B})) \end{array}
\]

并给出一个例子 *（在实直线上）*，其中这三个集合两两不同。

\begin{enumerate}
\def\labelenumi{\arabic{enumi})}
\setcounter{enumi}{5}
\item
  证明：拓扑空间 X 的子集 A 与 X 的每个稠密子集相交，当且仅当 A 的内部非空。
\item
  考虑给定拓扑空间 \(\mathbf{X}\) 的以下四个性质：
\end{enumerate}

(D \(_{I}\) ) X 的拓扑具有可数基。

\((D_{II})\) X 有一个可数稠密子集。

\((D_{III})\) X 的每个所有点都是孤立点的子集是可数的。

\((D_{IV})\) X 的两两不相交的非空开子集所成的每个集是可数的。

证明 \((\mathbf{D}_{\mathbf{I}})\) 蕴含 \((\mathbf{D}_{\mathbf{II}})\) 与 \((\mathbf{D}_{\mathbf{III}})\)，且 \((\mathbf{D}_{\mathbf{II}})\) 与 \((\mathbf{D}_{\mathbf{III}})\) 各自蕴含 \((\mathbf{D}_{\mathbf{IV}})\)（*）。

\begin{enumerate}
\def\labelenumi{\arabic{enumi})}
\setcounter{enumi}{7}
\item
  若拓扑空间 X 的子集 A 没有孤立点，则它在 X 中的闭包 \(\overline{A}\) 也没有孤立点。
\item
  设 X 是一个集合，\(M \to \overline{M}\) 是 \(\mathfrak{P}(X)\) 到其自身上的一个映射，使得：1) \(\overline{\varnothing} = \varnothing\)；2) 对每个 \(M \subset X\) 有 \(M \subset \overline{M}\)；3) 对每个 \(M \subset X\) 有 \(\overline{M} = \overline{M}\)；4) 对所有 \(M \subset X\) 与 \(N \subset X\) 有
\end{enumerate}

\[
\overline {{\mathbf {M} \cup \mathbf {N}}} = \overline {{\mathbf {M}}} \cup \overline {{\mathbf {N}}}.
\]

证明 X 上存在唯一的拓扑，使得 \(\overline{M}\)（对所有 \(M \subset X\)）是 M 关于这个拓扑的闭包（借助其闭集来定义这个拓扑）。

